

\documentclass[11pt]{article}
\usepackage[margin=1in]{geometry}
\usepackage{amsmath,amssymb,amsthm}
\usepackage{graphicx}
\usepackage{booktabs}
\usepackage{enumitem}
\usepackage{tabularx}
\usepackage{microtype}
\usepackage[hypertexnames=false]{hyperref}
\usepackage{algorithm}
\usepackage{algpseudocode}


\title{Notarized Anonymity Certificates Under Retries:\\
Grinding-Resistant Evaluation via Transparency Logs, Public Randomness, and \texorpdfstring{$\delta$}{delta}-Spending}
\author{}
\date{Draft (2026-03-17; archive v394)}

\newtheorem{theorem}{Theorem}
\newtheorem{lemma}{Lemma}
\newtheorem{proposition}{Proposition}

% Shared notation/style across the Anonymity / Anonymous DHT paper series.

\newcommand{\TV}{\mathrm{TV}}
\newcommand{\E}{\mathbb{E}}
\newcommand{\Pp}{\mathbb{P}}

% Reduce overfull boxes without forcing full \sloppy.
\setlength\emergencystretch{2em}

\begin{document}
\maketitle

\begin{abstract}
Empirical anonymity claims for anonymous overlays and DHTs are often published only after repeated evaluation, metric tuning, and parameter sweeps until ``the numbers look good.''
This stop-when-pass workflow can drive the probability of approving an unsafe build arbitrarily close to $1$, even when each individual check has false-approval probability at most $\delta$.
We give a compact evaluation discipline robust to retries: (i) a summable $\delta$-spending schedule across attempts (and within-attempt metric libraries), (ii) notarized manifests that precommit code, parameters, and evaluation randomness, and (iii) publication of every initiated attempt to an append-only transparency log with beacon/VRF binding so aborts and selective disclosure are auditable.
We prove that the disclosed budget upper-bounds false approvals under optional stopping and include toy simulations illustrating inflation under naive workflows and recovery under log-and-spend.
\end{abstract}

\paragraph{Compilation note.}
This paper is brief-by-design. A companion addendum collects extended figures, log-hardening discussion, and implementation templates (Eval~1A).\cite{series:evalnotaryaddendum}

\paragraph{Scope.}
We specify a \emph{notarized} evaluation workflow for anonymity claims: how to bind a reported privacy budget to a concrete claim object and downstream release, so that retries, audits, and later packaging remain checkable.
Within the now-staged retry / runtime / claim / release stack, this paper is the \emph{heavy evaluator-notarization companion}: it owns the attempt-log, many-testing, and verifier-report cut that makes later claims replayable without itself becoming a runtime-conformance or release-lineage note.

\paragraph{Imports.}
We import (i) schedule-only retry wrappers and downgrade-surface discipline (Operational~A), (ii) auditable trace interfaces (Operational~B; Synthesis~3), and (iii) ABOM/OINL claim binding (Anonymity~C). We treat the beacon/VRF binding and transparency-log interface as owned by the primitives note (Synthesis~9).\cite{series:transparencyprimitives}
Budget semantics are imported from odds inflation and the endpoint bridge (Anonymity~A; Synthesis~5).\cite{series:oddsinflation,series:endpointbridge,series:abomoinl}
Notation conventions for keys, state, transcripts, and transcript projections follow the series notation sheet (Synthesis~8). Exposure-normal-form identity (ENF-ID) is imported from the trace-interface note (Synthesis~3), and threat/window declarations (TW-ID) from the threat-window note (Synthesis~4); receipt-lineage semantics are imported from the receipt-schema note. Release-ledger comparison and PVSA packaging are imported from Release~A.\cite{series:notation,series:traceifaces,series:threatwindows,series:receiptitems,series:releaseproperty}
This paper owns only the \emph{attempt-log / spending / verifier} layer that sits below cross-release comparison.

\paragraph{Exports.}
Exports a minimal notarized evaluation bundle: a manifest schema, a gap-free logged attempt prefix, randomness binding, and a short verifier report (including many-testing adjustments) that reproduces the stated anonymity bound under the declared retry policy.
These objects can be referenced by ABOM/OINL and packaged by Release~A, but they do not themselves decide cross-release lineage or PVSA policy.\cite{series:abomoinl,series:releaseproperty}

\paragraph{Artifact policy.}
This archive is source-only; supporting artifacts (logs, traces, scripts, certificates, and verifier outputs) are available on request as a single bundle, committed by ABOM/OINL and governed by the series artifact policy (Synthesis~6).\cite{series:artifactpolicy}


\paragraph{Certificate object (informal).}
Exports a grinding-resistant \emph{notarized evaluation bundle}: public randomness, sample commitments, run digests, a gap-free logged attempt prefix,
and an auditor-verifiable trail that pins down retries and many-testing adjustments.
\paragraph{Contributions.}
\begin{itemize}
\item We formalize how retries, parameter sweeps, and selective disclosure can turn ``level-$\delta$'' evaluations into near-certain acceptance, even when a single run is conditionally valid.
\item We give a compact mitigation: a public \emph{spending schedule} across logged attempts, coupled with pre-commitment to beacon/VRF-derived randomness, so that aborts and retries cannot inflate false releases beyond the disclosed budget.
\item We provide a notarized publication checklist (manifest + transparency log + witness/gossip hooks) that lets readers verify the retry discipline without needing all artifacts inline.
\end{itemize}

\section{Context in the series}
This note is a \emph{deployment wrapper} around prior technical work in the series:
the privacy/anonymity objective is the same (e.g., TV leakage bounds, stopping-time leakage, and padding compilers),
and we aim to publish short certificates rather than full artifacts.
We therefore do \emph{not} re-derive the anonymity objective or padding theory; see
\cite{series:schedcompiler,series:spectral,series:knapsacktransport,series:manytestinglowerbounds}.
Instead, we address the practical question:
\emph{how do we publish a certificate that remains valid when the evaluator can retry, sweep, and selectively disclose?} A companion note develops advantage contracts for stealth audits and audit-aware behavior; see \cite{series:advantagecontracts}.

\paragraph{Series vocabulary (shared).}
We use a small, consistent vocabulary across the series to keep papers non-redundant.
A \emph{plan id} names an evaluation definition (trace schema, attempt unit, within-attempt library/correction, and spending schedule);
changing any evaluation knob implies a new plan id.
An \emph{epoch} (or window) is the declared validity interval of a plan (e.g., a wall-clock period or a beacon-derived id); changing epochs or restart conditions should be explicit.
An \emph{attempt} is one execution of a plan id over a predeclared epoch/window; retries, aborts, and detector changes are attempts and are logged.
A \emph{manifest} is a signed, hash-addressed statement committing code hash, plan id, parameters, and a randomness binding, and pointing to an offline artifact bundle (``available on request'').

\paragraph{Stable handoff to the release layer.}
This note intentionally stops at the earliest stage that answers the retry question.
Its owned audit spine is: \emph{plan id / manifest $\rightarrow$ gap-free logged attempt prefix $\rightarrow$ verifier report}.
Operational~A and Operational~B fix the retry wrapper and deployment-conformance surface beneath that spine; ABOM/OINL then names the resulting claim object and exposes the reader-facing diff surface; Release~A decides cross-release lineage, signed receipt issuance, and PVSA policy; and Eval~1A remains only the optional extended-figure / template companion.\cite{series:abomoinl,series:releaseproperty,series:evalnotaryaddendum}
A predecessor pointer inside an evaluation packet is therefore only a \emph{comparison hint}; by itself it does not settle public lineage.

\paragraph{Later-terse-paper citation posture.}
When a later short paper only needs to answer \emph{which replayable evaluation bundle made this anonymity number non-gameable under retries?}, it should stop at the minimal public notarized-evaluation card below plus the tiny spending drill.
It should widen only when the live issue has shifted to claim-surface comparison or release-lineage / PVSA policy, which remain owned by ABOM/OINL plus Release~A.\cite{series:abomoinl,series:releaseproperty}

\section{Problem: retries turn \texorpdfstring{``level-$\delta$''}{``level-delta''} into ``almost surely pass''}
Consider an evaluation attempt that outputs a statistic and a binary decision \emph{approve / reject}.
If the build is unsafe (violates the claimed bound), assume the attempt approves with probability at most $\delta$.
A common release gate repeats attempts until the first approval:
\[
\text{approve if } \exists t \le T \text{ such that attempt $t$ approves.}
\]
Under independence this approves an unsafe build with probability $1-(1-\delta)^T$, which can be close to $1$ even for modest $T$.
Worse, the evaluator can adaptively vary slices, metrics, or random seeds (``many testing'') \cite{series:manytestinglowerbounds}.

\paragraph{This is the security analogue of ``p-hacking''.}
The statistical failure mode here is not special to anonymity: it is the same ``researcher degrees of freedom'' problem as optional stopping and undisclosed analysis flexibility,
popularized as \emph{p-hacking} and the \emph{garden of forking paths} \cite{simmons2011falsepositive,gelman2013forkingpaths}.
In our setting the degrees of freedom are operational (which traces, which slices, which parameters, when to stop), but the consequence is identical:
nominal ``level-$\delta$'' gates can be driven toward near-certain approval unless retries and within-attempt sweeps are explicitly budgeted and logged.

\paragraph{Many metrics per attempt are also retries.}
Even if attempts are budgeted with a summable schedule, a common workflow is to run a \emph{library} of metrics/slices per attempt and publish the most favorable result.
Under an unsafe null, sweeping $M$ metrics at level $\alpha$ inflates the per-attempt false approval probability to $1-(1-\alpha)^M$.
Figure~\ref{fig:within_attempt_library} illustrates that attempt-level spending alone does not prevent this inflation; the per-attempt budget $\delta_t$ must also be split across the within-attempt library
(e.g., Bonferroni $\delta_t/M$, or another precommitted multiple-testing correction; see also \cite{series:manytestinglowerbounds}).

\begin{figure}[t]
\centering
\fbox{\begin{minipage}{0.92\linewidth}\centering\small Figure omitted in source-only release. Requested file: \texttt{figures/within\_attempt\_metric\_library.png}.\end{minipage}}
\caption{Within-attempt many-testing inflation: even with a summable spending schedule across attempts, sweeping many metrics/slices inside each attempt inflates false approvals unless the attempt budget is split (e.g., Bonferroni).}
\label{fig:within_attempt_library}
\end{figure}




\begin{table}[t]
\centering
\begin{tabularx}{\linewidth}{@{}lX@{}}
\toprule
Object & Practical budgeting rule \\
\midrule
Attempt index $t=1,2,\dots$ & Choose summable weights $w_t$ (e.g., $w_t \propto 1/t^2$ or $w_t \propto 2^{-t}$); set $\delta_t=\delta\, w_t$.\\
Within-attempt library & Precommit $M = |\mathcal{M}|\cdot|\mathcal{S}|\cdot|\mathcal{R}|$ (metrics $\times$ slices $\times$ seeds/restarts) and run each check at level $\delta_t/M$ (Bonferroni) or another precommitted correction \cite{series:manytestinglowerbounds}.\\
Aborts / retries & Treat each abort/retry as consuming its attempt budget $\delta_t$ (and log it), so grinding is explicitly paid for.\\
Publication & Publish only the first attempt whose full library passes; log the full manifest (including the library definition and all results).\\
\bottomrule
\end{tabularx}
\caption{A concrete $\delta$-budgeting recipe that blocks grinding across attempts and within-attempt sweeps with minimal operational overhead.}
\label{tab:budget_recipe}
\end{table}

Figure~\ref{fig:stop_when_pass} shows this inflation in a toy model.
The companion addendum includes an analogous illustration for metric/slice/seed sweeping and other within-attempt grinding patterns. \cite{series:evalnotaryaddendum}


\begin{figure}[t]
\centering
\fbox{\begin{minipage}{0.92\linewidth}\centering\small Figure omitted in source-only release. Requested file: \texttt{figures/release\_stop\_when\_pass.png}.\end{minipage}}
\caption{Naive stop-when-pass release gates inflate false safety approvals (unsafe releases).}
\label{fig:stop_when_pass}
\end{figure}


% (Extended material moved to Eval~1A addendum.)


\section{Threat model: how published anonymity certificates get gamed}
We focus on \emph{evaluation-aware} release workflows:
\begin{itemize}
\item \textbf{Stop-when-pass / rerun-until-pass.} Repeat the same check until a favorable draw.
\item \textbf{Selective disclosure.} Run many checks; publish only the successful ones and hide failures.
\item \textbf{Post-hoc commitment.} Commit to code/parameters after observing outcomes.
\item \textbf{Abort-after-beacon grinding.} If the evaluation randomness is derived from a public beacon, abort runs that look unfavorable.
\end{itemize}
We assume the evaluator can keep private artifacts; we only require that \emph{any published certificate} be backed by a verifiable manifest and log entries.

\section{Anytime-valid spending for repeated attempts}
Retries are a form of \emph{optional stopping}: if the evaluator can keep rerunning until something ``looks good,'' nominal level-$\delta$ gates cease to mean much.
The simplest fix is to precommit a \emph{spending schedule} $\{\delta_t\}_{t\ge 1}$ with $\sum_t\delta_t\le\delta_{\mathrm{tot}}$.
This is the time-Bonferroni analogue of the \emph{alpha-spending} approach in sequential analysis \cite{lan_demets1983}, and it matches the spirit of anytime-valid inference based on e-values/e-processes \cite{ville1939,shafer2021betting,howard2020timeuniform,ramdas2023gts_savi}.
The series also treats related ``stop-time'' pitfalls on the \emph{privacy side} (adaptive deadline choices) in a companion addendum \cite{series:stoptimeaddendum}.
Concretely, attempt $t$ outputs either (i) a p-value $P_t$ (valid under the unsafe null) or (ii) an e-value $E_t$ (with $\E[E_t]\le 1$ under the unsafe null).

\subsection{Summable spending rule}
Fix a sequence $\{\delta_t\}_{t\ge 1}$ with $\sum_t \delta_t \le \delta_{\text{tot}}$.
Approve at the first $t$ such that either $P_t \le \delta_t$ (p-value gate) or $E_t \ge 1/\delta_t$ (e-value gate).

\begin{lemma}[Precommitted spending is safe under retries]\label{lem:precommit}
Under the unsafe null, if each $P_t$ is a valid p-variable (or each $E_t$ is a valid e-variable),
then the spending gate approves with probability at most $\delta_{\text{tot}}$, even if the evaluator chooses a stopping time adaptively.
\end{lemma}

\begin{proof}[Proof sketch]
For p-variables, $\Pp(P_t \le \delta_t)\le \delta_t$ and a union bound yields $\Pp(\exists t: P_t\le \delta_t)\le \sum_t \delta_t$.
For e-variables, Markov's inequality gives $\Pp(E_t \ge 1/\delta_t)\le \delta_t$ and the same bound applies.
More refined proofs appear in the anytime-valid inference literature \cite{ramdas2023gts_savi}.
\end{proof}

\begin{lemma}[Spending with a within-attempt library]\label{lem:within_attempt}
Suppose attempt $t$ evaluates a precommitted library of $M_t$ metrics/slices and approves if \emph{any} metric passes.
If each metric test within attempt $t$ uses threshold $\delta_t/M_t$ (Bonferroni), then under the unsafe null the overall approval probability across all attempts is at most $\sum_t \delta_t \le \delta_{\mathrm{tot}}$ (even under adaptive stopping).
\end{lemma}

\begin{proof}[Proof sketch]
Under the unsafe null, within attempt $t$ the union bound gives
$\Pp(\text{attempt $t$ approves}) \le \sum_{m\le M_t} \delta_t/M_t = \delta_t$.
Applying Lemma~\ref{lem:precommit} to the attempt-level events yields the result.
\end{proof}



\subsection{Abort is spend}
If an evaluator can \emph{abort} after seeing information correlated with outcomes (e.g., a beacon-derived seed),
aborts must count as spent attempts.

\begin{lemma}[Abort-after-beacon grinding must be budgeted]\label{lem:abort}
If an evaluator can conditionally abort attempts after learning any partial information about the attempt outcome,
then a spending proof requires counting each initiated attempt (including aborted ones) against the spending schedule.
\end{lemma}

\begin{proof}[Proof sketch]
The spending guarantee (Lemma~\ref{lem:precommit}) bounds $\Pp(\exists t:\text{approve at }t)$ by $\sum_t\delta_t$ only when each initiated attempt is accounted for.
If aborts are free, the evaluator can implement rejection sampling: reveal a beacon-derived seed $U\sim\mathrm{Unif}[0,1]$, define a valid p-variable as $P=U$, and abort whenever $U>\delta_1$.
Since aborted attempts do not consume budget, the evaluator can repeat until $U\le\delta_1$ occurs, which happens almost surely, forcing approval with probability $1$ despite a nominal level-$\delta_1$ gate.
Counting every initiated attempt (including aborted ones) against the schedule restores the union-bound argument, because each peek/abort consumes the corresponding $\delta_t$.
\end{proof}


The companion addendum includes a toy illustration of abort-after-beacon grinding (and partial-leakage grinding) and how log-and-spend blocks it. \cite{series:evalnotaryaddendum}



% (Extended material moved to Eval~1A addendum.)


\section{Notarized manifests and transparency logs}
Spending controls retry inflation, but it does \emph{not} prevent \emph{selective disclosure}: an evaluator could run many attempts privately and publish only the favorable one.
To make the retry budget auditable, we require that every \emph{initiated} attempt be committed \emph{before} outcomes are observed and later accounted for in an append-only public record.

We use a standard transparency-log pattern: each attempt produces a signed, hash-addressed \emph{manifest} that is appended to an append-only Merkle log with inclusion and consistency proofs.
This is the same commitment surface used in Certificate Transparency (CT) \cite{rfc9162} and in supply-chain transparency logs such as Sigstore/Rekor \cite{newman2022sigstore}.
For concreteness, Sigstore bundles can embed transparency-log entry material---including a Signed Entry Timestamp (SET, an inclusion promise), an inclusion proof plus checkpoint, and optionally an RFC3161 timestamp---which together support offline verification of log inclusion and signing time.\cite{sigstore_bundle_spec,sigstore_threat_model}
The paper keeps artifacts off-net: the manifest carries only hashes and pointers; the full archive remains ``available on request.''

\subsection{Manifest}

Each attempt is associated with a signed \emph{manifest} containing:
\begin{itemize}
\item code hash (commit, container digest, or tarball hash),
\item evaluation plan identifier (metric family / slice library / spending schedule),
\item parameter set and environment fingerprint,
\item traffic model identifier (key distribution, probe window/epoch, vantage selection),
\item bootstrap / snapshot hash (routing-table snapshot, peer-set capture, or bootstrap list used to initialize the evaluation),
\item adversary / threat model identifier (e.g., compromise fraction and capability assumptions) when the evaluation depends on it,
\item randomness commitment: either a VRF output \cite{rfc9381,micali1999vrf} or a public beacon round (e.g., drand \cite{drand_spec} or NIST beacon \cite{nist_beacon,nist_ir8213_2019}); for robustness, it is reasonable to hash-mix multiple independent beacons when available,
\item the attempt index $t$ (which binds the spending threshold $\delta_t$).
\end{itemize}
Artifacts can remain off-chain; the manifest points to a retrievable archive by hash.

\paragraph{Portable encoding (optional).}
To keep certificates machine-checkable across ecosystems, a manifest can be encoded as an \emph{in-toto Statement} wrapped in a \emph{DSSE} envelope, matching common transparency ecosystems (e.g., Sigstore bundles) \cite{intoto_statement_spec,intoto_envelope_spec,dsse_spec,sigstore_bundle_spec}.
The protocol here only needs the \emph{hash-addressed} statement; full artifacts remain ``available on request.''
Recent IETF work on Supply Chain Integrity, Transparency, and Trust (SCITT) is standardizing a general architecture for signed-statement transparency services; our manifests can be viewed as SCITT signed statements and the log evidence as receipt material, enabling cross-ecosystem tooling reuse \cite{ietf_scitt_architecture}.

% (Extended material moved to Eval~1A addendum.)


\begin{table}[t]
\centering
\small
\begin{tabularx}{0.97\linewidth}{@{}p{0.27\linewidth}X@{}}
\toprule
Field & Meaning \\
\midrule
Subject hash & Code or build identifier (commit, container digest, or tarball hash). \\
Predicate type & For example \texttt{.../anonymity-certificate/v1}. \\
Plan id & Metric/slice library plus spending-schedule identifier. \\
Exposure normal-form id & Stable identifier or schema hash for the transcript projection whose budget is being certified. \\
Threat-window id & Stable identifier or schema hash for the repetition unit plus window, horizon, universe, and conditioning declaration (Synthesis~4). \\
Traffic model id & Key distribution plus probe window/epoch and vantage selection; binds evaluation degrees of freedom. \\
Bootstrap/snapshot hash & Hash of the bootstrap list, routing snapshot, or peer-set capture used to initialize the attempt. \\
Threat model id (optional) & Identifier for adversary assumptions the evaluation depends on, such as capabilities or compromise fraction. \\
Comparison hint (optional) & Prior certificate id plus relation tag used only to request later ABOM/ledger comparison; not sufficient by itself for public release lineage. \\
Attempt index $t$ & Binds the threshold $\delta_t$. \\
Randomness commitment & Precommitted VRF output or beacon round. \\
Artifact hash & Hash of the retrievable offline archive. \\
Log context & Log URL or id plus checkpoint reference. \\
\bottomrule
\end{tabularx}
\caption{Minimal notarized manifest fields (portable encoding).}
\label{tab:manifest_fields}
\end{table}

\paragraph{Distributing trust roots.}
To reduce reliance on out-of-band configuration, publish the log's URL(s), public key(s), witness policy, and rotation schedule using an update-security framework such as TUF \cite{tuf_spec_latest}.


\subsection{Append-only log}
Manifests are appended to an append-only log, yielding a publicly auditable record of all attempts.
A standard implementation is a Merkle-tree transparency log with consistency proofs (as in Certificate Transparency) \cite{rfc9162}.

If the publication surface is simply ``signed checksums plus proofs'' (rather than an arbitrary statement API), minimalist transparency designs such as Sigsum are also suitable and come with monitor-oriented tooling and offline verification workflows \cite{sigsum_site,sigsum_design,sigsum_getting_started}.

In practice, modern deployments increasingly use tile-backed, read-optimized logs (the Static CT API) to make monitoring cheaper and more robust at scale (e.g., Sunlight and other tile-based logs built atop Tessera) \cite{ct_static_2025,sunlight_intro_2024,static_ct_api_spec,tessera_github}.
Mirroring protocols make it easier to keep independent read-only copies of the log (and to serve tiles from CDNs/object storage), improving availability for monitors and referees. \cite{c2sp_tlog_mirror}
For portability, the log inclusion and consistency evidence can be packaged as an offline-verifiable proof bundle (e.g., C2SP \texttt{tlog-proof} / \texttt{tlog-checkpoint} formats) so referees can verify a certificate without contacting the log operator \cite{c2sp_tlog_proof,c2sp_tlog_checkpoint}.
Figure~\ref{fig:log_diagram} shows the minimal data flow.

\begin{figure}[t]
\centering
\fbox{\begin{minipage}{0.92\linewidth}\centering\small Figure omitted in source-only release. Requested file: \texttt{figures/release\_transparency\_log.png}.\end{minipage}}
\caption{Publishing notarized manifests to a transparency log binds retries, spending, and randomness.}
\label{fig:log_diagram}
\end{figure}

\subsection{Protocol}
\begin{algorithm}[t]
\caption{Log-enforced attempt}\label{alg:log_attempt}
\begin{algorithmic}[1]
\Require spending thresholds $\{\delta_t\}$, attempt index $t$, beacon/VRF randomness source
\State create manifest $M_t$ with code hash, plan id, exposure normal-form id, parameters, beacon round / VRF output, and $t$
\State append $M_t$ to transparency log; obtain inclusion proof
\State derive evaluation seed from (manifest, beacon/VRF) and run evaluation
\State publish result (statistic + p/e-value) with log inclusion proof
\end{algorithmic}
\end{algorithm}

\paragraph{Deterministic seed derivation.}
A concrete default is
$\mathrm{seed}_t := H(\,H(M_t)\,\|\,\mathrm{beacon}(r_t)\,\|\,t\,)$,
where $H$ is a fixed hash (e.g., SHA-256), $H(M_t)$ is the manifest hash, and $\mathrm{beacon}(r_t)$ is the public randomness value for the committed round $r_t$ (or a VRF output).
This makes the evaluation randomness \emph{recomputable} and blocks post-hoc seed choice.

\paragraph{Canonical hashing of structured manifests.}
If $M_t$ is a structured object (e.g., JSON), then ``hashing the manifest'' must specify a canonical serialization;
otherwise, semantically identical manifests can hash differently across languages or tooling.
Two safe defaults are: (i) hash the raw bytes of a signed envelope payload (e.g., DSSE), or (ii) use a standard JSON canonicalization scheme such as RFC~8785 (JCS) before hashing.\cite{rfc8785}



\paragraph{Gap-free attempt indices and commit-charged spending.}
To make retries auditable, attempt indices should be gap-free within a plan id and epoch: the log (or a monitor) flags missing indices as aborted attempts.
Operationally, charge the attempt budget at \emph{commit time} (when the manifest is appended) rather than at completion time, and treat missing reveals/results as spent aborts (or an operational alarm), so abort-after-beacon grinding cannot concentrate budget on favorable instances.

\begin{proposition}[When two notarized evaluation certificates are locally comparable]
\label{prop:notary-lineage}
Suppose a publisher wants to compare a newly notarized evaluation certificate against a prior one and report the result as a local evaluation delta.
This is meaningful only if the new certificate declares either
(i) the same plan id, exposure normal form id, repetition unit, and public conditioning as the old one, or
(ii) an explicit coarsening of the old exposure normal form under the same plan-level repetition unit and public conditioning.
Otherwise the correct action is to mint a new claim object (with an optional predecessor pointer) and hand comparison off to ABOM/OINL plus the release ledger, rather than to treat the change as a simple update inside one continuing evaluation surface.
\end{proposition}

\begin{proof}[Proof sketch]
The retry-spending guarantee says how to account for \emph{repeated attempts of one declared experiment}.
If the plan id or covered transcript projection changes, the experiment itself has changed, so the statement ``the later certificate is better/worse than the earlier one'' is no longer a comparison of like with like.
Explicit coarsening is the one safe exception because it preserves the old statement up to data processing on the public view.
Public release lineage is a stronger claim and belongs to the ABOM/receipt layer above this proposition.
\end{proof}

\paragraph{Earliest sufficient stop.}
If the question is ``was this evaluation run under a non-gameable retry discipline?'' an audit may stop once the spending schedule, the gap-free logged attempt prefix, and the verifier report all check out.
If the question is instead ``how does this compare to the previous public claim?'' the audit must hand off to ABOM/OINL plus the release ledger before reading any headline delta.\cite{series:abomoinl,series:releaseproperty}

\begin{table}[t]
\centering
\small
\begin{tabularx}{0.97\linewidth}{@{}p{0.18\linewidth}p{0.43\linewidth}X@{}}
\toprule
\textbf{Field} & \textbf{Minimal public value} & \textbf{Why this row is enough at the notarization stop} \\
\midrule
Plan anchor & \texttt{plan\_id = eval.notary.retry.v1}; spending schedule \texttt{sched.delta.4step.v1}; max attempts $T_{\max}=4$ & Pins the exact repeated-attempt game being audited \\
Claim tuple & \texttt{exposure\_nf\_id = enf.retry.trace.v1}; \texttt{tw\_id = tw.lookup.30d.v1}; repetition unit = one declared epoch & Prevents ``same certificate, different public view'' drift \\
Logged prefix & Attempts $1,2$ present; attempt $2$ is declared as the first passing attempt & Makes hidden favorable restarts inside the disclosed prefix auditable \\
Randomness binding & Manifest digest plus beacon round \texttt{drand:312345} determine the evaluation seed & Blocks post-hoc seed grinding \\
Many-testing rule & Attempt $2$ used a five-metric library with Bonferroni threshold $\delta_2/5$ & Shows the within-attempt correction rather than only the headline pass bit \\
Verifier verdict & $p_2 = 4\times 10^{-8}\le \delta_2/5 = 6\times 10^{-8}$; no earlier attempt crossed its threshold & Replays the acceptance arithmetic in one line \\
Transparency anchor & Inclusion proofs for attempts $1,2$ checked at checkpoint \texttt{cp-912}; consistency verified from prior checkpoint & Makes the logged prefix publicly discoverable and append-only checkable \\
Artifact commitment & Replay bundle digest \texttt{sha256:7c31\ldots a8e0} committed; full archive available on request & Keeps the public card small without making the private bundle unauditable \\
\bottomrule
\end{tabularx}
\caption{A minimal public notarized-evaluation card. This is the smallest public answer later terse papers can quote directly when the live issue is whether one anonymity number came from a replayable, non-gameable evaluation bundle under retries.}
\label{tab:minimal-notary-card}
\end{table}

Table~\ref{tab:minimal-notary-card} is intentionally non-decorative: it names one plan anchor, one logged attempt prefix, one randomness binding, one correction rule, and one verifier verdict, without silently widening into runtime-conformance receipts or release-lineage packets.

\paragraph{One tiny spending drill.}
Suppose the disclosed spending schedule is $(\delta_1,\delta_2,\delta_3,\delta_4)=(4,3,2,1)\times 10^{-7}$, so the total published budget is $10^{-6}$.
If attempt $2$ is the first passing attempt and its within-attempt metric library has size $M_2=5$, then Bonferroni allows acceptance only when a reported metric satisfies
\[
p \le \delta_2/M_2 = 3\times 10^{-7}/5 = 6\times 10^{-8}.
\]
A later terse paper may therefore say ``the headline number came from logged attempt $2$ under a disclosed $10^{-6}$ retry budget'' without reopening the full evaluator archive.

\subsection{Selective disclosure: what the log buys you (brief)}
Spending controls retry inflation, but does not by itself prevent an evaluator from \emph{hiding} failed attempts.
The point of logging is that, once every initiated attempt is appended \emph{before} evaluation (Algorithm~\ref{alg:log_attempt}),
the evaluator cannot later publish only the favorable outcome without leaving a public record of the failed (or aborted) attempts.

A referee-facing acceptance rule is therefore simple:
accept a certificate only if it includes (i) log inclusion proofs for a gap-free prefix of attempt indices $t=1,\dots,t^\star$ under the stated plan id, exposure normal-form id, and epoch,
(ii) a globally consistent checkpoint for those inclusions (e.g., witnessed/cosigned), and
(iii) evidence that attempt $t^\star$ is the \emph{first} attempt in that prefix to cross its threshold.
Under this rule, \emph{selective disclosure inside the protocol} becomes auditable: missing indices indicate aborts (still spent), and omitted manifests cannot be hidden because they are already in the public log.

\subsection{Split-view resistance: witnesses, gossip, and cross-logging (brief)}
If the log operator can equivocate (show different histories to different verifiers), then log inclusion proofs cease to be a global commitment.
Standard defenses are (i) checkpoint gossip/monitoring, (ii) witnessed checkpoints (cosigned tree heads), and optionally (iii) cross-logging checkpoints into an independent log.
These mechanisms are orthogonal to the retry-spending guarantee, but they are needed if third parties are to treat the log as authoritative; see \cite{dahlberg2018ctgossip,syta2016cosi,dirksen2021logpicker,hicks2023sok_transparency}. Recent standards work is also beginning to cite these witness/cosignature building blocks (e.g., Merkle Tree Certificates) \cite{ietf_plants_merkle_tree_certs}.

Operationally, this is becoming easier: public witness networks and directories reduce bespoke bootstrap and make it realistic for referees to demand witnessed checkpoints (or equivalent gossip evidence) for published claims \cite{witness_network_blog_2024,trustfabric_witnesses_2025,filippo_keyserver_tlog_2025}.


% (Extended material moved to Eval~1A addendum.)

\begin{algorithm}[t]
\caption{Notarized evaluation gate (approve-once)}\label{alg:release_gate}
\begin{algorithmic}[1]
\Require total budget $\delta_{\text{tot}}$, summable schedule $\{\delta_t\}$, max attempts $T_{\max}$
\For{$t=1$ to $T_{\max}$}
  \State run Algorithm~\ref{alg:log_attempt} and obtain p-value $P_t$ (or e-value $E_t$)
  \If{$P_t \le \delta_t$ \textbf{or} $E_t \ge 1/\delta_t$}
     \State \textbf{publish} evaluation bundle: manifest list, log proofs, and the accepted attempt index $t$
     \State \Return
  \EndIf
\EndFor
\State \textbf{publish} failure bundle (optional): manifest list and log proofs
\end{algorithmic}
\end{algorithm}
\subsection{Reader verification checklist (minimal, artifact-light)}
A published certificate can stay \emph{small} while still being referee-checkable.
At minimum, publish (or link to) a compact bundle containing: the spending schedule identifier, the accepted attempt index $t^\star$,
the manifest(s) and log proofs for attempts $1,\dots,t^\star$, the beacon/VRF evidence, and the accepted test statistic.

A referee (or reader) can then verify, mechanically:
\begin{enumerate}
\item \textbf{Budget accounting.} Confirm the disclosed schedule satisfies $\sum_t \delta_t \le \delta_{\text{tot}}$ and that each initiated attempt consumes its $\delta_t$
(including aborts); confirm any within-attempt library uses a precommitted correction (e.g., Bonferroni $\delta_t/M_t$).
\item \textbf{Precommitment + claim surface.} For each attempt $t\le t^\star$, verify that the signed manifest (code hash, plan id, exposure normal-form id, parameters, traffic model, snapshot hash) was \emph{logged} before outcomes,
using the inclusion proof and an append-only checkpoint/consistency proof.
\item \textbf{Global log view.} Verify the checkpoint is globally consistent (e.g., via a witnessed/cosigned checkpoint and published witness policy, or via monitor/gossip evidence) \cite{c2sp_tlog_witness,c2sp_tlog_cosignature}.
\item \textbf{Randomness binding.} Verify the evaluation seed is deterministically derived from the manifest and the stated beacon round or VRF output, so the evaluator could not grind by post-hoc seed choice.
\item \textbf{Gate correctness.} Verify the published $P_{t^\star}$ (or $E_{t^\star}$) crosses the disclosed threshold $\delta_{t^\star}$ (or $1/\delta_{t^\star}$), and that the certificate is indeed the \emph{first} passing attempt in the log.
\item \textbf{Local comparability discipline.} If the certificate is presented as comparable to an earlier evaluation output, verify the declared comparison hint is compatible with Proposition~\ref{prop:notary-lineage}; otherwise treat the output as a new claim object and defer any public lineage claim to ABOM/OINL plus Release~A.
\end{enumerate}
This is enough to make the retry discipline auditable \emph{without} publishing full artifacts inline; hashes and ``available on request'' archives suffice.



% (Extended material moved to Eval~1A addendum.)



\section{Practical checklist for publishing anonymity certificates}
\noindent Table~\ref{tab:dht_degrees_freedom} lists common ``degrees of freedom that silently change trace distributions in anonymous-DHT evaluations;
for a non-gameable certificate, bind them into the plan id / manifest rather than leaving them as undocumented evaluator choices.

\begin{table}[t]
\centering
\begin{tabularx}{\linewidth}{@{}lX@{}}
\toprule
Degree of freedom & What to commit (minimum) \\
\midrule
Key distribution & Uniform hashes vs prefix-targeted / adversarial keys; key batch size; any reserved ranges. \\
Vantage selection & How start nodes are sampled; IP/AS diversity policy; whether probes come from stable monitors vs churny clients. \\
Lookup semantics & Iterative vs recursive; parallelism $\alpha$; termination rule; timeout / retry rule; caching behavior. \\
Routing state & Bootstrap list / snapshot hash; routing-table freshness policy; churn model (if simulated). \\
Overlay parameters & $k$ (bucket size), ID length, eviction policy, replication factor; any protocol extensions (e.g., BEP-5 details \cite{bep5_mainline_dht}, libp2p/IPFS Kad-DHT specifics \cite{libp2p_kad_dht_spec,ipfs_kad_dht_spec}). \\
Measurement transport & UDP/TCP choices, pacing, connection reuse; any middlebox traversal; clock sync assumptions. \\
\bottomrule
\end{tabularx}
\caption{Anonymous-DHT evaluation knobs that should be fixed \emph{up front} (plan id / manifest) to avoid hidden degrees of freedom.}
\label{tab:dht_degrees_freedom}
\end{table}

\paragraph{Example: protocol semantics are part of the plan.}
In practice, ``Kademlia-style'' DHTs vary in which messages propagate and who answers lookups.
For example, I2P's network database uses a floodfill subset for netDb queries: store messages are flooded among floodfills using a Kademlia XOR metric, but leaseSet lookup queries are answered by the receiving floodfill without forwarding (to avoid a security issue).
This changes both the set of in-path observers and the detectability surface for audits, so such forwarding/answering rules belong in the committed attempt definition.\cite{i2p_tech_intro}

A certificate intended to be non-gameable should include:
\begin{enumerate}
\item Manifest schema, including code hash, parameters, the evaluation plan id, and the exposure normal-form id of the transcript projection being certified.
\item An explicit \emph{attempt unit} definition for your anonymous DHT evaluation: the number of lookups per attempt, how keys are sampled (uniform, prefix-targeted, etc.), how start nodes/vantage points are sampled, and the epoch/window over which traces are collected. Bind this definition into the plan id and/or manifest (e.g., for Kademlia-family lookups \cite{maymounkov2002kademlia}).
\item A fixed spending schedule $\{\delta_t\}$ and an explicit statement of $\delta_{\text{tot}}$. A convenient default is $\delta_t \propto 1/t^2$ (summable with a heavy head), or a geometric schedule when early approval is expected. If each attempt sweeps a library of metrics/slices, split $\delta_t$ across that library (e.g., Bonferroni $\delta_t/M_t$) as in Lemma~\ref{lem:within_attempt}.
\item A public log inclusion proof for every initiated attempt (including aborted ones).
\item A randomness binding: VRF output \cite{rfc9381,micali1999vrf} or a public beacon \cite{drand_spec,nist_beacon} fixed before outcomes are observed.
\item A short ``artifacts on request'' statement keyed by archive hashes.
\end{enumerate}

\section{Takeaways}
\begin{itemize}[leftmargin=*]
\item \textbf{Retries are a release attack:} repeated evaluation and selective disclosure can inflate false approvals to near certainty even if each run is ``level-$\delta$''.
\item \textbf{Log-and-spend is the minimal fix:} bind every attempt (including aborts) into an append-only log, and publish a summable $\delta$-spending schedule (plus within-attempt splits) so optional stopping cannot exceed the disclosed budget.
\item \textbf{Bind randomness and degrees of freedom:} commit probe selection via VRF/beacon and treat evaluation semantics (keys, vantages, retries, timeouts, caching) as plan-id parameters.
\item \textbf{Make verification cheap:} a short manifest + logged prefix + verifier report lets third parties audit the retry discipline without shipping full artifacts inline, while leaving lineage and signed release packaging to the release layer.
\end{itemize}

\section{Limitations and open problems}
\paragraph{Dependence and drift.}
Spending controls Type-I error under optional stopping, but the underlying evaluation statistic must remain valid in the presence of dependence and distribution shift.
In anonymous DHTs, dependence arises from repeated peers, caching, and diurnal latency.
A conservative operational answer is to treat major configuration changes as new experiments (restart the spending schedule),
and/or to run block-wise tests with a mixing assumption; tighter integration with overlay-specific dependence models is left for future work.

\paragraph{Log equivocation.}
If an operator can maintain multiple inconsistent logs, transparency loses force.
Witnessing and gossip reduce the feasibility of split-view attacks but do not remove the need for diverse operational assumptions (independent witnesses/monitors, failure modes, and on-path adversaries).
When high assurance is required, require witness cosigning (or cross-logging checkpoints to independent logs) and treat witness-set changes as configuration changes that restart certification.



% (Extended material moved to Eval~1A addendum.)


\begin{thebibliography}{99}\small

\bibitem{maymounkov2002kademlia}
Petar Maymounkov and David Mazières.
\newblock Kademlia: A peer-to-peer information system based on the XOR metric.
\newblock In \emph{IPTPS}, 2002.

\bibitem{bep5_mainline_dht}
BitTorrent Enhancement Proposal 5.
\newblock DHT Protocol (Mainline).
\newblock \url{https://www.bittorrent.org/beps/bep_0005.html}.

\bibitem{libp2p_kad_dht_spec}
libp2p.
\newblock Kademlia DHT specification.
\newblock \url{https://github.com/libp2p/specs/blob/master/kad-dht/README.md}.

\bibitem{ipfs_kad_dht_spec}
IPFS.
\newblock IPFS Kademlia DHT specification.
\newblock \url{https://specs.ipfs.tech/routing/kad-dht/}.

\bibitem{drand_spec}
drand.
\newblock Protocol specification.
\newblock \url{https://docs.drand.love/docs/specification/}.

\bibitem{nist_beacon}
NIST.
\newblock Interoperable Randomness Beacons / NIST Randomness Beacon.
\newblock \url{https://csrc.nist.gov/projects/interoperable-randomness-beacons}.

\bibitem{micali1999vrf}
Silvio Micali, Michael Rabin, and Salil Vadhan.
\newblock Verifiable Random Functions.
\newblock In \emph{FOCS}, 1999.

\bibitem{syta2016cosi}
Ewa Syta et al.
\newblock Keeping Authorities ``Honest or Bust'' with Decentralized Witness Cosigning.
\newblock In \emph{IEEE S\&P}, 2016.

\bibitem{rfc9162}
Certificate Transparency v2.
\newblock RFC 9162.
\newblock \url{https://www.rfc-editor.org/rfc/rfc9162.html}.

\bibitem{ct_static_2025}
C2SP.
\newblock Static Certificate Transparency API.
\newblock \url{https://github.com/C2SP/C2SP/blob/main/static-ct-api.md}.

\bibitem{cloudflare_azul_ct_2025}
Chris Hartwig and Vincent Van Doordrecht.
\newblock A next-generation Certificate Transparency log built on Cloudflare Workers.
\newblock Cloudflare blog, 2025.
\newblock \url{https://blog.cloudflare.com/azul-certificate-transparency-log/}.

\bibitem{dirksen2021logpicker}
Alexandra Dirksen et al.
\newblock LogPicker: Strengthening Certificate Transparency Against Covert Adversaries.
\newblock \emph{PoPETs}, 2021.
\newblock \url{https://petsymposium.org/popets/2021/popets-2021-0066.pdf}.

\bibitem{dahlberg2018ctgossip}
Rasmus Dahlberg et al.
\newblock Aggregation-Based Certificate Transparency Gossip.
\newblock \texttt{arXiv:1806.08817}, 2018.

\bibitem{hicks2023sok_transparency}
Alexander Hicks.
\newblock SoK: Log Based Transparency Enhancing Technologies.
\newblock \texttt{arXiv:2305.01378}, 2023.

\bibitem{rfc8785}
JSON Canonicalization Scheme.
\newblock RFC 8785.
\newblock \url{https://www.rfc-editor.org/rfc/rfc8785}.

\bibitem{rfc9381}
Verifiable Random Functions (VRFs).
\newblock RFC 9381.
\newblock \url{https://www.rfc-editor.org/rfc/rfc9381.html}.

\bibitem{series:stoptimeaddendum}
Anonymous.
\newblock Stop-Time Padding Addendum (unique-only): Max-Mixture Separations and Tail-Sign Deadline Normal Forms.
\newblock Series note (published), 2026.
\newblock Wiki: \texttt{[[2026.01.26 - Mathematics: Stop-Time Padding Addendum, Max-Mixture Separations and Tail-Sign Deadline Normal Forms]]}.

\bibitem{newman2022sigstore}
Zachary Newman et al.
\newblock Sigstore: Software Signing for Everybody.
\newblock In \emph{CCS}, 2022.



\bibitem{series:transparencyprimitives}
Anonymous.
\newblock Transparency Primitives for Proof-Carrying Anonymity Claims: Beacons, VRFs, and Append-Only Logs.
\newblock Series synthesis note (Synthesis~9), The-One-True-Archive (2026).

\bibitem{series:artifactpolicy}
Anonymous.
\newblock Artifact and Disclosure Policy for the One-True-Archive (Source-Only Public Release).
\newblock Series synthesis note (Synthesis~6), 2026. (This archive.)


\bibitem{series:notation}
Anonymous.
\newblock Notation and Minimal Model for the One-True-Archive.
\newblock Series synthesis note (Synthesis~8), 2026. (This archive.)

\bibitem{series:traceifaces}
Anonymous.
\newblock Trace interfaces for anonymous DHTs.
\newblock Synthesis~3, 2026. (This archive.)

\bibitem{series:threatwindows}
Anonymous.
\newblock Threat models and repetition windows for anonymous DHTs.
\newblock Synthesis~4, 2026. (This archive.)


\bibitem{series:receiptitems}
Anonymous.
\newblock Receipt line items: a minimal tuple schema for privacy claims.
\newblock Synthesis~12, 2026. (This archive.)

\bibitem{series:oddsinflation}
Anonymous.
\newblock Odds Inflation, Privacy Loss, and Endpoint Anonymity Bounds.
\newblock Anonymity series Paper~A, 2026. (This archive.)

\bibitem{series:endpointbridge}
Anonymous.
\newblock Endpoint Metrics Bridge: Odds Inflation, PML/MaxL, and Identification Sizing.
\newblock Series synthesis note (Synthesis~5), 2026. (This archive.)

\bibitem{series:abomoinl}
Anonymous.
\newblock ABOM and OINL for Anonymity Claims.
\newblock Anonymity series Paper~C, 2026. (This archive.)

\bibitem{series:releaseproperty}
Anonymous.
\newblock Privacy as a Release Property: Proof-Carrying Anonymity Receipts for Anonymous DHTs.
\newblock Release and destination Paper~A, 2026. (This archive.)

\bibitem{series:schedcompiler}
Anonymous.
\newblock \emph{Scheduling and Compiler Techniques for Metadata-Hiding Overlay Lookups}.
\newblock Series paper (published), 2026.
\newblock Wiki: \texttt{[[2026.01.22 - Mathematics: Scheduling and Compiler Techniques for Metadata-Hiding Overlay Lookups]]}.

\bibitem{series:spectral}
Anonymous.
\newblock \emph{Spectral Anonymity and Optimal Distinguishing Bounds for Random-Walk Delegation in (Possibly Directed) Overlays}.
\newblock Series paper (published), 2026.
\newblock Wiki: \texttt{[[2026.01.23 - Mathematics: Spectral Anonymity and Optimal Distinguishing Bounds for Random-Walk Delegation in (Possibly Directed) Overlays]]}.

\bibitem{series:knapsacktransport}
Anonymous.
\newblock \emph{Optimal Poset-Feasible Padding: Knapsack, Transport, and Dual Certificates for TV Privacy}.
\newblock Series paper (published), 2026.
\newblock Wiki: \texttt{[[2026.01.29 - Mathematics: Optimal Poset-Feasible Padding: Knapsack, Transport, and Dual Certificates for TV Privacy]]}.

\bibitem{series:manytestinglowerbounds}
Anonymous.
\newblock \emph{Many-Testing Lower Bounds for Low-Latency Privacy on Hierarchical Overlays}.
\newblock Series paper (published), 2026.

\bibitem{series:advantagecontracts}
\newblock Advantage Contracts for Anonymous DHT Evaluations: Stealth Audits, Leakage Discovery at Scale, and Practical Padding Filters.
\newblock Online manuscript in the author's Anonymity / Anonymous DHT series, 2026.

\bibitem{sunlight_intro_2024}
{Let's Encrypt}.
\newblock Introducing Sunlight: a Certificate Transparency log implementation built for the Static CT API.
\newblock Let's Encrypt Blog, 2024.
\newblock Accessed 2026-02
\newblock \url{https://letsencrypt.org/2024/03/14/introducing-sunlight}

\bibitem{sunlight_reflections_2025}
{Let's Encrypt}.
\newblock Reflections on a Year of Sunlight.
\newblock Let's Encrypt Blog, 2025.
\newblock Accessed 2026-02
\newblock \url{https://letsencrypt.org/2025/06/11/reflections-on-a-year-of-sunlight}

\bibitem{sunlight_site}
\newblock The Sunlight CT Log (Static CT API).
\newblock sunlight.dev, 2026.
\newblock Accessed 2026-02
\newblock \url{https://sunlight.dev/}

\bibitem{static_ct_api_spec}
\newblock Static Certificate Transparency API.
\newblock C2SP specification, 2026.
\newblock Spec for tile/static CT endpoints (also mirrored at c2sp.org); accessed 2026-02
\newblock \url{https://c2sp.org/static-ct-api}

\bibitem{itko_blog_2025}
Easterday, Adit.
\newblock I Built a New Certificate Transparency Log in 2024 -- Here's What I Learned.
\newblock transparency.dev blog, 2025.
\newblock Accessed 2026-02
\newblock \url{https://blog.transparency.dev/i-built-a-new-certificate-transparency-log-in-2024-heres-what-i-learned}

\bibitem{chrome_ct_policy_2024}
\newblock What Chrome is thinking about Certificate Transparency in 2024 and beyond.
\newblock Chromium CT Policy mailing list, 2024.
\newblock Accessed 2026-02
\newblock \url{https://groups.google.com/a/chromium.org/g/ct-policy/c/W7OSO3SbrFo}

\bibitem{chrome_ct_log_policy_2026}
\newblock Certificate Transparency Log Policy.
\newblock Chrome Certificate Transparency documentation, 2026.
\newblock States acceptance of {RFC 6962} logs and logs conforming to the {static-ct-api} {C2SP} specification; accessed 2026-02
\newblock \url{https://googlechrome.github.io/CertificateTransparency/log_policy.html}

\bibitem{apple_ct_program_2026}
\newblock Apple's Certificate Transparency log program.
\newblock Apple Support, 2026.
\newblock Mentions {static-ct-api} conformance for logs; accessed 2026-02
\newblock \url{https://support.apple.com/en-us/103703}

\bibitem{rekor_v2_ga_2025}
\newblock Rekor v2 GA -- Cheaper to run, simpler to maintain.
\newblock Sigstore Blog, 2025.
\newblock Accessed 2026-02
\newblock \url{https://blog.sigstore.dev/rekor-v2-ga/}

\bibitem{openssf_rekor_dataset_2025}
{OpenSSF}.
\newblock Announcing the Sigstore Transparency Log Research Dataset.
\newblock OpenSSF Blog, 2025.
\newblock Accessed 2026-02
\newblock \url{https://openssf.org/blog/2025/10/15/announcing-the-sigstore-transparency-log-research-dataset/}

\bibitem{trailofbits_rekor_monitor_2025}
{Trail of Bits}.
\newblock Catching malicious package releases using a transparency log.
\newblock Trail of Bits Blog, 2025.
\newblock Accessed 2026-02
\newblock \url{https://blog.trailofbits.com/2025/12/12/catching-malicious-package-releases-using-a-transparency-log/}

\bibitem{nist_ir8213_2019}
Kelsey, John.
\newblock {NIST IR 8213}: A Reference for Randomness Beacons (Format and Protocol Version 2).
\newblock NIST CSRC, 2019.
\newblock Initial public draft published May 2019; accessed 2026-02
\newblock \url{https://csrc.nist.gov/pubs/ir/8213/ipd}

\bibitem{c2sp_tlog_tiles}
\newblock Transparency Log Tiles API (tlog-tiles).
\newblock C2SP specification, 2026.
\newblock Accessed 2026-02
\newblock \url{https://c2sp.org/tlog-tiles}

\bibitem{c2sp_tlog_mirror}
\newblock Transparency Log Mirrors (tlog-mirror).
\newblock C2SP specification, 2026.
\newblock HTTP protocol for mirroring transparency logs; accessed 2026-02
\newblock \url{https://c2sp.org/tlog-mirror}

\bibitem{c2sp_tlog_checkpoint}
\newblock Transparency Log Checkpoint Format (tlog-checkpoint).
\newblock C2SP specification, 2026.
\newblock Accessed 2026-02
\newblock \url{https://c2sp.org/tlog-checkpoint}

\bibitem{c2sp_tlog_proof}
\newblock Transparency Log Offline Proof Format (tlog-proof).
\newblock C2SP specification, 2026.
\newblock Accessed 2026-02
\newblock \url{https://c2sp.org/tlog-proof}

\bibitem{c2sp_tlog_witness}
\newblock Transparency Log Witness Protocol (tlog-witness).
\newblock C2SP specification, 2026.
\newblock Accessed 2026-02
\newblock \url{https://c2sp.org/tlog-witness}

\bibitem{c2sp_tlog_cosignature}
\newblock Transparency Log Cosignature Format (tlog-cosignature).
\newblock C2SP specification, 2026.
\newblock Accessed 2026-02
\newblock \url{https://c2sp.org/tlog-cosignature}

\bibitem{tessera_github}
\newblock Tessera: A Go library for tile-based transparency logs.
\newblock transparency-dev GitHub repository, 2026.
\newblock Accessed 2026-02
\newblock \url{https://github.com/transparency-dev/tessera}

\bibitem{tesseract_blog_2025}
\newblock Introducing TesseraCT.
\newblock blog.transparency.dev, 2025.
\newblock Announcement of TesseraCT (Static CT API implementation) built atop tile-backed transparency logs; accessed 2026-02
\newblock \url{https://blog.transparency.dev/introducing-tesseract}

\bibitem{trillian_site}
\newblock Trillian: General Transparency (project site).
\newblock google.github.io/trillian, 2026.
\newblock Notes that Trillian is in maintenance mode and recommends tile-backed logs via Tessera; accessed 2026-02
\newblock \url{https://google.github.io/trillian/}

\bibitem{filippo_keyserver_tlog_2025}
Valsorda, Filippo.
\newblock Building a Transparent Keyserver.
\newblock words.filippo.io, 2025.
\newblock Practical integration of tile-backed tlogs; accessed 2026-02
\newblock \url{https://words.filippo.io/keyserver-tlog/}

\bibitem{sigsum_site}
\newblock Sigsum: Key-usage transparency with signed checksums.
\newblock sigsum.org, 2026.
\newblock Minimalist transparency log ecosystem; accessed 2026-02
\newblock \url{https://www.sigsum.org/}

\bibitem{sigsum_design}
\newblock Sigsum design document.
\newblock sigsum project documentation, 2026.
\newblock Design overview, including log, witness, and monitor roles; accessed 2026-02
\newblock \url{https://git.sigsum.org/sigsum/tree/doc/design.md}

\bibitem{sigsum_getting_started}
\newblock Sigsum: Getting started (offline verification demo).
\newblock sigsum.org documentation, 2026.
\newblock Step-by-step workflow: log a checksum and verify offline; accessed 2026-02
\newblock \url{https://www.sigsum.org/getting-started/}

\bibitem{tesseract_repo}
\newblock Tesseract: An implementation of the Static CT API based on Tessera.
\newblock transparency-dev GitHub repository, 2026.
\newblock Accessed 2026-02
\newblock \url{https://github.com/transparency-dev/tesseract}

\bibitem{tuf_spec_latest}
\newblock The Update Framework (TUF) Specification (Latest).
\newblock theupdateframework.github.io, 2026.
\newblock Accessed 2026-02
\newblock \url{https://theupdateframework.github.io/specification/latest/}

\bibitem{dsse_spec}
\newblock DSSE: Dead Simple Signing Envelope (Specification).
\newblock secure-systems-lab GitHub repository, 2026.
\newblock Accessed 2026-02
\newblock \url{https://github.com/secure-systems-lab/dsse}

\bibitem{intoto_envelope_spec}
\newblock in-toto Attestation Framework v1: Envelope layer (DSSE).
\newblock in-toto/attestation specification, 2026.
\newblock Accessed 2026-02
\newblock \url{https://github.com/in-toto/attestation/blob/main/spec/v1/envelope.md}

\bibitem{intoto_statement_spec}
\newblock in-toto Attestation Framework v1: Statement layer.
\newblock in-toto/attestation specification, 2026.
\newblock Accessed 2026-02
\newblock \url{https://github.com/in-toto/attestation/blob/main/spec/v1/statement.md}

\bibitem{slsa_provenance_v1}
\newblock SLSA Provenance v1.0.
\newblock slsa.dev specification, 2026.
\newblock Accessed 2026-02
\newblock \url{https://slsa.dev/spec/v1.0/provenance}

\bibitem{sigstore_bundle_spec}
\newblock Sigstore Bundle Format.
\newblock sigstore.dev documentation, 2025.
\newblock Accessed 2026-02
\newblock \url{https://docs.sigstore.dev/about/bundle/}

\bibitem{trillian_tessera_blog_2024}
\newblock Introducing Trillian Tessera: Tile-based Transparency Logs.
\newblock blog.transparency.dev, 2024.
\newblock Accessed 2026-02
\newblock \url{https://blog.transparency.dev/introducing-trillian-tessera}

\bibitem{tessera_v1_blog_2025}
\newblock Tessera v1.0 is here.
\newblock blog.transparency.dev, 2025.
\newblock Accessed 2026-02
\newblock \url{https://blog.transparency.dev/tessera-v1-0-is-here}

\bibitem{witness_network_blog_2024}
\newblock Can I Get a Witness (Network)?.
\newblock blog.transparency.dev, 2024.
\newblock Overview of public witness networks; accessed 2026-02
\newblock \url{https://blog.transparency.dev/can-i-get-a-witness-network}

\bibitem{trustfabric_witnesses_2025}
\newblock TrustFabric Witnesses (conforming to C2SP tlog-witness).
\newblock transparency.dev, 2025.
\newblock List of public/dev witnesses; accessed 2026-02
\newblock \url{https://transparency.dev/witnesses/}

\bibitem{simmons2011falsepositive}
Simmons, Joseph P. and Nelson, Leif D. and Simonsohn, Uri.
\newblock False-Positive Psychology: Undisclosed Flexibility in Data Collection and Analysis Allows Presenting Anything as Significant.
\newblock Psychological Science, 22(11), 1359--1366, 2011.
\newblock DOI: 10.1177/0956797611417632.
\newblock Classic discussion of inflated false positives under undisclosed analysis flexibility

\bibitem{gelman2013forkingpaths}
Gelman, Andrew and Loken, Eric.
\newblock The Garden of Forking Paths: Why Multiple Comparisons Can Be a Problem, Even When There Is No ``Fishing Expedition'' or ``p-hacking'' and the Research Hypothesis Was Posited Ahead of Time.
\newblock Unpublished manuscript, 2013.
\newblock Accessed 2026-02
\newblock \url{https://sites.stat.columbia.edu/gelman/research/unpublished/p_hacking.pdf}

\bibitem{lan_demets1983}
Lan, K. K. Gordon and DeMets, David L..
\newblock Discrete sequential boundaries for clinical trials.
\newblock Biometrika, 70(3), 659--663, 1983.
\newblock DOI: 10.1093/biomet/70.3.659.

\bibitem{ville1939}
Ville, Jean.
\newblock \'{E}tude critique de la notion de collectif.
\newblock 1939.
\newblock Th\`ese; accessed 2026-02
\newblock \url{https://www.numdam.org/item/THESE_1939__218__1_0.pdf}

\bibitem{howard2020timeuniform}
Howard, Steven R. and Ramdas, Aaditya and McAuliffe, Jon and Sekhon, Jasjeet.
\newblock Time-uniform Chernoff bounds via nonnegative supermartingales.
\newblock Probability Surveys, 17, 257--317, 2020.
\newblock DOI: 10.1214/18-PS321.
\newblock \url{https://projecteuclid.org/journals/probability-surveys/volume-17/issue-none/Time-uniform-Chernoff-bounds-via-nonnegative-supermartingales/10.1214/18-PS321.full}


\bibitem{ramdas2023gts_savi}
Ramdas, Aaditya and Gr\"unwald, Peter and Vovk, Vladimir and Shafer, Glenn.
\newblock Game-theoretic statistics and safe anytime-valid inference.
\newblock \emph{Statistical Science}, 38(4), 2023.
\newblock \texttt{arXiv:2210.01948}.

\bibitem{shafer2021betting}
Shafer, Glenn.
\newblock Testing by betting: A strategy for statistical and scientific communication.
\newblock Journal of the Royal Statistical Society: Series A, 184(2), 407--431, 2021.
\newblock DOI: 10.1111/rssa.12647.
\newblock \url{https://rss.onlinelibrary.wiley.com/doi/abs/10.1111/rssa.12647}

\bibitem{melara2015coniks}
Melara, Marcela S. and Blankstein, Aaron and Bonneau, Joseph and Felten, Edward W. and Freedman, Michael J..
\newblock CONIKS: Bringing Key Transparency to End Users.
\newblock 24th USENIX Security Symposium (USENIX Security 15), 2015.
\newblock \url{https://www.usenix.org/system/files/conference/usenixsecurity15/sec15-paper-melara.pdf}

\bibitem{sun2024ctwatchers}
Sun, Aozhuo and Lin, Jingqiang and Wang, Wei and Liu, Zeyan and Li, Bingyu and Wen, Shushang and Wang, Qiongxiao and Li, Fengjun.
\newblock Certificate Transparency Revisited: The Public Inspections on Third-party Monitors.
\newblock Network and Distributed System Security (NDSS) Symposium, 2024.
\newblock DOI: 10.14722/ndss.2024.24834.
\newblock Proposes \"watchers\" to inspect third-party CT monitors and detect inconsistencies
\newblock \url{https://www.ndss-symposium.org/wp-content/uploads/2024-834-paper.pdf}

\bibitem{tomescu2019transparencylogs}
Tomescu, Alin and Bhupatiraju, Vivek and Papadopoulos, Dimitrios and Papamanthou, Charalampos and Triandopoulos, Nikos and Devadas, Srinivas.
\newblock Transparency Logs via Append-Only Authenticated Dictionaries.
\newblock ACM SIGSAC Conference on Computer and Communications Security (CCS), 2019.
\newblock CCS '19; transparency logs and equivocation/monitoring tradeoffs; accessed 2026-02
\newblock \url{https://www.cs.yale.edu/homes/cpap/published/logs-ccs19.pdf}

\bibitem{ietf_plants_merkle_tree_certs}
\newblock Merkle Tree Certificates.
\newblock IETF Internet-Draft (draft-ietf-plants-merkle-tree-certs), 2026.
\newblock Internet-Draft; current versions cite C2SP tlog specs (e.g., tlog-checkpoint/tlog-mirror/tlog-tiles); accessed 2026-02
\newblock \url{https://datatracker.ietf.org/doc/draft-ietf-plants-merkle-tree-certs/}

\bibitem{i2p_tech_intro}
\newblock I2P Technical Introduction: Network Database (floodfill / netDb).
\newblock I2P Project documentation, 2026.
\newblock See ``Network Database'' section (floodfill, store propagation, and lookup behavior); accessed 2026-02
\newblock \url{https://geti2p.net/ca/docs/how/tech-intro}

\bibitem{i2p_netdb_overview}
\newblock The Network Database (netDb) and Floodfill.
\newblock I2P Project documentation, 2026.
\newblock Overview of floodfill roles, routing keys, and lookup messages; accessed 2026-02
\newblock \url{https://i2p.net/en/docs/overview/network-database/}

\bibitem{ietf_scitt_architecture}
\newblock An Architecture for Trustworthy and Transparent Digital Supply Chains.
\newblock IETF Internet-Draft: draft-ietf-scitt-architecture, 2025.
\newblock SCITT architecture for signed-statement transparency services; accessed 2026-02
\newblock \url{https://datatracker.ietf.org/doc/html/draft-ietf-scitt-architecture-22}

\bibitem{rekor_overview}
\newblock Rekor Transparency Log: Logging Overview.
\newblock Sigstore documentation, 2026.
\newblock Describes Rekor entries, inclusion proofs, and verification tooling; accessed 2026-02
\newblock \url{https://docs.sigstore.dev/logging/overview/}

\bibitem{sigstore_threat_model}
\newblock Sigstore Threat Model.
\newblock Sigstore documentation, 2026.
\newblock Describes the Sigstore workflow and Rekor inclusion promises (SET) in the threat model; accessed 2026-02
\newblock \url{https://docs.sigstore.dev/about/threat-model/}

\bibitem{series:evalnotaryaddendum}
Anonymous.
\newblock Notarized Anonymity Certificates: Extended Figures and Log-Hardening Notes (Addendum).
\newblock Evaluation series companion note (Eval~1A), 2026. (This archive.)

\end{thebibliography}

\end{document}